% Abhakseite „Das kann ich" – erzeugt von abhaken_bau.py aus den Quelltexten
\begin{abhakseite}
\abhakgruppe{Kennst du schon}
\abhak{1}{Ich kann Zahlen quadrieren, auch negative Zahlen und Dezimalzahlen.}
\abhak{2}{Ich kann mit Klammern und negativen Zahlen rechnen.}
\abhak{3}{Ich kann durch Einsetzen prüfen, ob eine Zahl eine Gleichung erfüllt.}
\abhak{4}{Ich finde den Fehler beim Einsetzen einer negativen Zahl.}
\abhak{5}{Ich kann eine negative Zahl in einen Term mit Quadrat einsetzen.}
\abhak{6}{Ich kann eine lineare Gleichung lösen.}
\abhak{7}{Ich kann eine Quadratwurzel ziehen.}
\abhak{8}{Ich kann den Scheitelpunkt aus der Scheitelpunktform ablesen.}
\abhak{9}{Ich kann Klammern ausmultiplizieren, auch mit binomischen Formeln.}
\abhak{10}{Ich kann $x$ ausklammern.}
\abhak{11}{Ich kann Terme ordnen und zusammenfassen.}
\abhakgruppe{Einheit 1 · Wurzelziehen und Lösbarkeit}
\abhak{12}{Ich erkenne, ob eine Gleichung quadratisch ist.}
\abhak{13}{Ich erkenne an der Zahl rechts, wie viele Lösungen eine Gleichung wie $x^2 = 5$ hat.}
\abhakgruppe{– Gleichungen der Form $x^2 = $ Zahl}
\abhak{14}{Ich kann eine Gleichung der Form $x^2 = $ Zahl durch Wurzelziehen lösen.}
\abhak{15}{Ich kann erst $x^2$ freistellen und dann die Wurzel ziehen.}
\abhakgruppe{– Gleichungen mit quadrierter Klammer}
\abhak{16}{Ich kann eine Gleichung mit quadrierter Klammer durch Rückwärtsrechnen lösen.}
\abhakgruppe{– Lösungen prüfen und zählen}
\abhak{17}{Ich kann prüfen, ob eine Zahl eine Lösung einer quadratischen Gleichung ist.}
\abhak{18}{Ich kann am Graphen ablesen, wie viele Lösungen eine Gleichung hat.}
\abhak{19}{Ich kann eine Gleichung angeben, die keine Lösung hat.}
\abhak{20}{Ich kann Aussagen über die Lösungen einer Gleichung beurteilen.}
\abhak{21}{Ich finde den Fehler beim Wurzelziehen.}
\abhak{22}{Ich kann begründen, warum eine Gleichung zwei, eine oder keine Lösung hat.}
\abhak{23}{Ich kann eine Länge durch Wurzelziehen berechnen.}
\abhakgruppe{Einheit 2 · Satz vom Nullprodukt}
\abhak{24}{Ich erkenne, ob ein Produkt gleich null gesetzt ist.}
\abhakgruppe{– Gleichungen in Produktform}
\abhak{25}{Ich kann eine Gleichung in Produktform mit dem Satz vom Nullprodukt lösen.}
\abhak{26}{Ich kann prüfen, ob eine Zahl eine Gleichung in Produktform erfüllt.}
\abhakgruppe{– Gleichungen, in denen jedes Glied ein $x$ hat}
\abhak{27}{Ich kann eine Gleichung, in der jedes Glied ein $x$ hat, durch Ausklammern lösen.}
\abhakgruppe{– Zwischen Produkt und Normalform wechseln}
\abhak{28}{Ich kann eine Produktgleichung mit einer Zahl rechts ausmultiplizieren und ordnen.}
\abhak{29}{Ich kann eine Gleichung lösen, wenn ich ihre Produktform kenne.}
\abhak{30}{Ich finde den Fehler beim Satz vom Nullprodukt.}
\abhak{31}{Ich kann Regeln zum Satz vom Nullprodukt begründen.}
\abhakgruppe{Einheit 3 · Normalform und p-q-Formel}
\abhak{32}{Ich erkenne, welche Form eine Gleichung hat und welcher Lösungsweg passt.}
\abhak{33}{Ich kann p und q mit Vorzeichen ablesen.}
\abhak{34}{Ich erkenne, ob die Vorzahl von $x^2$ eins ist.}
\abhakgruppe{– Lösen mit der p-q-Formel}
\abhak{35}{Ich kann eine Gleichung in Normalform mit der p-q-Formel lösen.}
\abhak{36}{Ich kann eine Gleichung erst ordnen oder normieren und dann mit der Formel lösen.}
\abhak{37}{Ich kann mit der Formel erkennen, ob eine Gleichung zwei, eine oder keine Lösung hat.}
\abhak{38}{Ich kann Lösungen mit Wurzel und als Näherungswert angeben.}
\abhak{39}{Ich kann mit einer Lösung die Probe machen.}
\abhakgruppe{– Gleichungen mit Klammern}
\abhak{40}{Ich kann eine Gleichung mit Klammern erst auflösen, dann ordnen und lösen.}
\abhakgruppe{– Den Lösungsweg wählen}
\abhak{41}{Ich kann eine Gleichung auf dem passenden Weg lösen.}
\abhak{42}{Ich finde den Fehler bei der p-q-Formel.}
\abhak{43}{Ich kann Regeln der p-q-Formel begründen.}
\abhak{44}{Ich kann die Schnittpunkte einer Geraden mit einer Parabel berechnen.}
\abhakgruppe{Einheit 4 · Sachaufgaben}
\abhak{45}{Ich erkenne, welche Lösung zur Frage passt.}
\abhakgruppe{– Zahlenrätsel}
\abhak{46}{Ich kann ein Zahlenrätsel mit einer quadratischen Gleichung lösen.}
\abhakgruppe{– Rechteckaufgaben}
\abhak{47}{Ich kann eine Rechteckaufgabe mit einer quadratischen Gleichung lösen.}
\abhak{48}{Ich finde den Fehler bei einer Sachaufgabe.}
\abhak{49}{Ich kann begründen, wann eine Sachaufgabe eine oder zwei Antworten hat.}
\end{abhakseite}
